\cleardoublepage

\section{总结与展望}

\subsection{本文工作总结}
本文在已有 ICESat-2 多次散射雪深反演理论基础上，面向 ATL03 光子数据构建了完整的实际处理与反演流程。围绕北极海冰区域样本，本文完成了 ATL03 强束光子读取、研究区筛选、主表面识别、表面对齐、后向散射剖面构建、背景噪声扣除和系统响应去卷积，并在统一处理框架下实现了基于路径长度统计矩的雪深迭代反演。

在理论依据方面，本文参考了 Hu 等关于多次散射路径长度分布和低阶统计矩关系的推导，将平均路径长度、扩散散射系数和雪深之间的物理联系作为反演基础\citep{Hu2022}。在实际处理方面，本文参考了 Lu 等关于 ICESat-2 雪深反演不确定性的分析，将 afterpulse 校正、表面粗糙度、背景噪声和空间配对误差作为影响反演结果的主要因素\citep{Lu2022Unc}。同时，Lu 等关于 ATLAS 海洋激光衰减系数的研究说明，ICESat-2 次表层光子剖面在不同介质中均具有光学信息提取潜力，但必须依赖严格的质量控制和剖面解释边界\citep{Lu2023OceanAtten}。这些文献共同构成了本文方法设计、结果解释和局限性讨论的主要参考。

在结果验证方面，本文选取 2019 年 4 月三期 ICESat-2 与 OIB 重叠观测样本进行空间配对分析，对反演结果与机载雪深进行定量比较。结果表明，本文方法能够在厘米量级误差范围内恢复样本区的背景雪深变化，整体 Bias 为 0.0058\,m，RMSE 为 0.0558\,m，MAE 为 0.0429\,m。与此同时，结果也表明该方法在深雪区仍存在振幅压缩，对局地峰谷变化的响应相对偏弱，说明剖面构建、去卷积稳定性和空间配对尺度仍会影响最终反演表现，相关误差需要结合观测尺度解释。本文还使用 2020 年 3 月 27 日 UA/CMC 格网雪深产品开展补充对比，结果显示 ICESat-2 反演雪深相对 UA 整体偏高，Bias 为 0.3253\,m，RMSE 为 0.9066\,m，相对 CMC 的 Bias 和 RMSE 分别为 0.5850\,m 和 0.9791\,m，表明当前海冰样本参数体系不能直接推广到陆地区域格网验证场景。

\subsection{后续展望}
本文工作表明，基于 ICESat-2 多次散射回波剖面的雪深直接反演在实际 ATL03 数据中具有可行性，但仍有进一步改进空间。后续研究可从以下几个方面展开。首先，可扩大样本区域和观测时段，分析不同冰型、不同地表类型、不同背景噪声条件和不同季节下参数设置的稳定性。当前样本主要集中在少数北极海冰重叠航段，尚不能充分代表多年冰、融化期雪层、陆地粗糙地形和深雪山区等复杂场景，因此需要在更多典型环境中重新检验路径长度矩与雪深之间的稳定关系。

其次，可进一步优化表面识别、背景扣除和系统响应校正方法，提高浅层长尾信号提取质量，并增强深雪区的振幅保持能力。具体而言，可结合 ATL03 提供的背景计数、太阳高度角、云和气溶胶信息，对昼夜观测和不同大气条件分别建立质量控制策略；也可借鉴海洋衰减系数反演中对表面回波强度、大气光学厚度和有效深度范围的筛选思想，对雪层剖面建立更严格的有效样本判据。对于去卷积结果，可进一步引入正则化和不确定性传播分析，以减少病态求解对深层统计矩的影响。

最后，可引入更多外部验证数据，对 OIB、UA、CMC 和现场观测等多源数据建立分场景验证口径，对反演结果开展更细致的多尺度比较与误差分析。不同验证数据代表的空间尺度和物理含义并不相同，OIB 更适合检验海冰沿轨局地雪深，UA/CMC 更适合检验陆地格网背景雪深，而现场观测可用于校准小尺度粗糙度和雪层结构影响。只有在验证尺度、地表类型和反演参数相互匹配的条件下，才能更可靠地评估该方法在大范围雪深监测中的适用性。
